\bumppoints{22}
\question\TFQuestion{T}{If $p=0.4$, then total homozygosity in the population is 0.52 (assume equilibrium).}
\question\TFQuestion{T}{Adaptations can lead to descent with modification.}
\question\TFQuestion{F}{One assumption of the Hardy-Weinberg equations is that mutations are random.}
\question\TFQuestion{T}{Complex organisms can have fewer chromosomes than simple organisms. }
\question\TFQuestion{T}{Genetic drift is strongest in small populations.}
\question\TFQuestion{T}{If a trait is energetically costly and provides no survival advantage, then natural selection may cause that trait to disappear from the population over many generations.}
\question\TFQuestion{F}{The genotype is determined by the interaction of the phenotype with the environment.}
\question\TFQuestion{F}{Microevolution states that individuals in a population evolve over time}.
\question\TFQuestion{T}{Gene flow between two populations can reduce the effects of genetic drift in each population.}
\question\TFQuestion{F}{Overproduction of offspring by every individual in a population is necessary for evolution to occur.}
\question\TFQuestion{T}{The Hardy-Weinberg equations show that genotype frequencies do not change if allele frequencies do not change.}
